\uuid{nYKP}
\titre{Matrice/Puissance/Suites}
\auteur{hanani ; mzari}
\organisation{exo7}
\licence{CC BY-NC-SA 4.0}
\datecreate{2020-01-20}
\niveau{L1}
\module{Algèbre}
\chapitre{Matrice}
\sousChapitre{Propriétés élémentaires, généralités}
\difficulte{3}
\theme{Calcul matriciel, Matrice, Puissance, Suites}
\exo7qcmid{lille-2.matrices.018}

\contenu{
\begin{question}
\qclassification{108.01}
\qtags{motcle=Matrice/Puissance/Suites}
Soit $A= \left(\begin{array}{rcc}1&2&3\\0&1&2\\ 0&0&1\\
\end{array}\right)$ et  $N= A-I$, où $I= \left(\begin{array}{rcc}
1&0&0\\0&1&0\\ 0&0&1\\ \end{array}\right).$\\
On considère $3$ suites récurrentes  $(u_n)_{n\ge 0}$, $(v_n)_{n\ge 0}$ et $(w_n)_{n\ge 0}$ définies par $u_0,v_0,w_0$ des réels donnés et pour $n\ge 1$ :
$$(\mathtt{S})\left\{\begin{array}{rcc}
u_n&=&u_{n-1}+2v_{n-1}+3w_{n-1}\\
v_n&=&v_{n-1}+2w_{n-1}\\ w_n&=&w_{n-1}. \\
\end{array}\right.$$
Quelles sont les assertions vraies ?
\begin{answers}
\bad{$N^k = 0$, pour tout entier $k\ge 2$}.
\good{Pour tout entier $n \ge 2,$ $A^n =I+nN+\frac{n(n-1)}{2}N^2  $}.
\bad{Pour tout entier $n \ge 0,$
$$(\mathtt{S}) \left\{\begin{array}{rcc}
u_n&=&u_0+2nv_0+3nw_0\\v_n&=&v_0+2nw_0\\ w_n&=&w_0. \\
\end{array}\right.$$}
\good{Pour tout entier $n \ge 0,$
$$(\mathtt{S})  \left\{\begin{array}{rcc}
u_n&=&u_0+2nv_0+n(2n+1)w_0\\v_n&=&v_0+2nw_0\\ w_n&=&w_0. \\
\end{array}\right.$$}
\end{answers}
\begin{explanations} On a : $N=\left(\begin{array}{rcc}
0&2&3\\0&0&2\\ 0&0&0\\
\end{array}\right)$, $N^2=\left(\begin{array}{rcc}
0&0&4\\0&0&0\\
0&0&0\\ \end{array}\right)$ et $N^k=0, $ pour tout $k\ge 3$.\\
Comme $A= N+I$ et le produit des matrices $N$ et $I$ est commutatif, on peut appliquer la formule du binôme pour le calcul des
puissances de $A$.
\vskip0mm
En calculant $A^n$ et  utilisant l'égalité : $\left(\begin{array}{r}
u_n\\v_n\\ w_n\\
\end{array}\right) = A^n \left(\begin{array}{r}
u_0\\v_0\\ w_0\\
\end{array}\right)$, on déduit que : $$\left\{\begin{array}{rcc}
u_n&=&u_0+2nv_0+n(2n+1)w_0\\
v_n&=&v_0+2nw_0\\
w_n&=&w_0. \\
\end{array}\right.$$
\end{explanations}
\end{question}
}
